\section{The \texttt{Mirror} McStas Component}
Single mirror plate.

\subsection*{Identification}
\begin{itemize}
  \item \textbf{Site:} 
  \item \textbf{Author:} Kristian Nielsen
  \item \textbf{Origin:} Risoe
  \item \textbf{Date:} August 1998
\end{itemize}

\subsection*{Description}
\begin{lstlisting}
Single mirror plate (used to build guides in compound components).
The component Mirror models a single rectangular neutron mirror plate.
It can be used as a sample component or to e.g. assemble a complete neutron
guide by putting multiple mirror components at appropriate locations and
orientations in the instrument definition, much like a real guide is build
from individual mirrors.
Additionally, it may be centered in order to be used as a sample or
monochromator plate, else it starts from AT position.
Reflectivity is either defined by an analytical model (see Component Manual)
or from a two-columns file with q [Angs-1] as first and R [0-1] as second.
In the local coordinate system, the mirror lies in the first quadrant of the
x-y plane, with one corner at (0, 0, 0). Plane is thus perpendicular to the
incoming beam, and should usually be rotated.

Example: Mirror(xwidth=.1, yheight=.1,R0=0.99,Qc=0.021,alpha=6.1,m=2,W=0.003)
\end{lstlisting}

\subsection*{Input parameters}
Parameters in \textbf{boldface} are required; the others are optional.

\begin{longtable}{p{0.22\textwidth}p{0.12\textwidth}p{0.46\textwidth}p{0.14\textwidth}}
\toprule
\textbf{Name} & \textbf{Unit} & \textbf{Description} & \textbf{Default} \\
\midrule
\endhead
reflect & str & Name of reflectivity file. Format q(Angs-1) R(0-1) & 0 \\
\textbf{xwidth} & m & width of mirror plate &  \\
\textbf{yheight} & m & height of mirror plate &  \\
R0 & 1 & Low-angle reflectivity & 0.99 \\
Qc & AA-1 & Critical scattering vector & 0.021 \\
alpha & AA & Slope of reflectivity & 6.07 \\
m & 1 & m-value of material. Zero means completely absorbing. & 2 \\
W & AA-1 & Width of supermirror cut-off & 0.003 \\
center & 1 & if true (1), the Mirror is centered arount AT position. & 0 \\
transmit & 1 & When true, non reflected neutrons are transmitted through the mirror, instead of being absorbed. & 0 \\
\bottomrule
\end{longtable}

\subsection*{Links}
\begin{itemize}
  \item \href{run:/home/willend/willend-McCode/mcstas-comps/optics/Mirror.comp}{Source code} for \texttt{Mirror.comp}.
\end{itemize}
\IfFileExists{Mirror_static.tex}{\input{Mirror_static.tex}}{}